%%%%%%%%%%%%%%%%%%%%%%%%%%%%%%%%%%%%%%%%%%%%%%%%%
% EXERCICE
\begin{exercise}[points=3]
    Complète le tableau ci-dessous.
    \begin{center}
        \begin{tblr}{
                hlines,
                vlines,
                columns={c,1cm},
                rows={c,1cm},
                column{5}={c,4cm},
                row{1}={font=\bfseries},
            }
            D        & d        & q        & r        & Égalité euclidienne \\
            55       & 5        &  &  &                     \\
             & 12       & 5        & 3        &                     \\
             &  &  &  & $29=3\cdot 8+5$               \\
        \end{tblr}
    \end{center}
\end{exercise}
% SOLUTION
\begin{solution}

\end{solution}
%%%%%%%%%%%%%%%%%%%%%%%%%%%%%%%%%%%%%%%%%%%%%%%%%
%%%%%%%%%%%%%%%%%%%%%%%%%%%%%%%%%%%%%%%%%%%%%%%%%
% EXERCICE
\begin{exercise}[points=6]
    Calcule le PGCD et le PPCM des nombres suivants.
    \begin{enumerate}[cols=3]
        \item 90 et 84.
        \item 15 et 30.
        \item 7 et 12.
    \end{enumerate}
    \vspace*{7cm}
\end{exercise}
% SOLUTION
\begin{solution}

\end{solution}
%%%%%%%%%%%%%%%%%%%%%%%%%%%%%%%%%%%%%%%%%%%%%%%%%
%%%%%%%%%%%%%%%%%%%%%%%%%%%%%%%%%%%%%%%%%%%%%%%%%
% EXERCICE
\begin{exercise}[points=2]
    Rends les fractions irréductibles.
    \begin{enumerate}[cols=2]
        \item $\dfrac{-144}{-108}=$
        \item $-\dfrac{120}{504}=$
    \end{enumerate}
    \vspace*{2cm}
\end{exercise}
% SOLUTION
\begin{solution}

\end{solution}
%%%%%%%%%%%%%%%%%%%%%%%%%%%%%%%%%%%%%%%%%%%%%%%%%
\newpage
%%%%%%%%%%%%%%%%%%%%%%%%%%%%%%%%%%%%%%%%%%%%%%%%%
% EXERCICE
\begin{exercise}[points=2]
Place sur la droite graduée suivante les points dont voici les abscisses. \\
    \hfil
    $abs A =\dfrac{3}{2}$
    \hfil
    $abs B =2$
    \hfil
    $abs C =-\dfrac{2}{5}$
    \hfil
    $abs D =-\dfrac{11}{10}$
    \hfil
    \begin{center}
        \begin{tikzpicture}[x=2.5cm, >=latex]
    \draw[->] (-3,0) -- (3,0)node[right]{$x$};
    \draw (0,0) node[below=3pt]{$0$};
    \draw (1,0) node[below=3pt]{$1$};
    \foreach \x in {-3,-2,...,2}
        {
            \draw (\x, -4pt) -- (\x, 4pt);
        }
    \foreach \x in {-3,-2.9,...,2.9}
        {
            \draw (\x, -2pt) -- (\x, 2pt);
        }
\end{tikzpicture}
    \end{center}
\end{exercise}
% SOLUTION
\begin{solution}

\end{solution}
%%%%%%%%%%%%%%%%%%%%%%%%%%%%%%%%%%%%%%%%%%%%%%%%%
% EXERCICE
\begin{exercise}[points=2]
    Encadre les nombres suivants au rang demandé.
    \begin{enumerate}[cols=2]
        \item $\dfrac{8}{7} \approx \numprint{1,142857}$ au centième près.
        \item $-\pi \approx \numprint{-3,141592}$ au dixième près.
    \end{enumerate}
    \vspace*{2cm}
\end{exercise}
% SOLUTION
\begin{solution}

\end{solution}
%%%%%%%%%%%%%%%%%%%%%%%%%%%%%%%%%%%%%%%%%%%%%%%%%
%%%%%%%%%%%%%%%%%%%%%%%%%%%%%%%%%%%%%%%%%%%%%%%%%
% EXERCICE
\begin{exercise}[points=6]
    Effectue au maximum en utilisant les propriétés des puissances.
    \begin{enumerate}[cols=2, itemsep=3cm]
        \item $\left(2ab\right)^3=$ %1
        \item $-3a^3\cdot a^4\cdot \left(y^3\right)^4=$ %2
        \item $\dfrac{6a^7}{a^3}=$ %1
        \item $\left( \dfrac{-2a}{3}\right)^3=$ %2
        \item $\left(-2a^2b\right)^4=$ %2    
        %\item $5x^2 \cdot \left(-x\right)^2 \cdot \left(3xy^4\right)^2$ %3
        \item $\left(a^3 \cdot a^5\right)^3= $ %2
        %\item $\left( \dfrac{-5x^4}{2x^2}\right)^3 $ %3
    \end{enumerate}
    \vspace*{3cm}
\end{exercise}
% SOLUTION
\begin{solution}

\end{solution}
%%%%%%%%%%%%%%%%%%%%%%%%%%%%%%%%%%%%%%%%%%%%%%%%%
\newpage
%%%%%%%%%%%%%%%%%%%%%%%%%%%%%%%%%%%%%%%%%%%%%%%%%
% EXERCICE
\begin{exercise}[points=2]
    À l'aide d'expressions algébriques, prouve les propriétés suivantes.
    \begin{enumerate}[itemsep=4cm]
        \item La somme de deux nombres consécutifs est toujours un nombre impair.
        \item La somme de deux nombres impairs est toujours un nombre pair.
    \end{enumerate}
    \vspace*{4cm}
\end{exercise}
% SOLUTION
\begin{solution}

\end{solution}
%%%%%%%%%%%%%%%%%%%%%%%%%%%%%%%%%%%%%%%%%%%%%%%%%